Discrete choice models provide a statistical language for understanding decisions among finite alternatives.
They are routinely used to study travel mode choice, route choice, product demand, pricing, and service operations, where researchers need both interpretable utility parameters and reliable counterfactual predictions \citep{mcfadden1974conditional,train2009discrete}.
As empirical choice data become larger and model structures become more heterogeneous, the computational requirements of choice modeling increasingly resemble those of modern differentiable systems: likelihoods must be vectorized, simulation draws must be reproducible, and gradients, Hessians, predictions, and post-estimation quantities must remain aligned.

Existing software serves important parts of this workflow, but no single open package currently combines PyTorch-native computation with econometric model abstractions and a systematic estimator benchmark suite.
Biogeme provides a mature Python environment for discrete choice model specification and estimation \citep{bierlaire2023biogeme}; Apollo provides a flexible R framework for a broad range of choice models \citep{hess2019apollo}; and R packages such as \texttt{mlogit} and \texttt{gmnl} provide widely used estimators for random utility models \citep{croissant2020mlogit,sarrias2017gmnl}.
These packages remain essential reference implementations, yet their computational backends and interfaces are not designed around PyTorch tensors, automatic differentiation pipelines, or direct integration with modern machine-learning workflows.
Conversely, PyTorch provides efficient tensor computation and automatic differentiation \citep{paszke2019pytorch}, but it does not itself provide the data structures, model syntax, inference outputs, and benchmark discipline expected in discrete choice research.

TorchDCM addresses this gap by treating software design and estimator validation as a single contribution.
The package keeps the model implementation in a pure Python/PyTorch library while placing heavyweight validation, reference-estimator wrappers, and public benchmark reports in a separate validation layer.
This separation is deliberate: the user-facing package should remain lightweight and installable, while the benchmark system can depend on Biogeme, Apollo, R packages, and large public datasets.
The result is a package that is both usable as a modeling library and auditable as an estimator implementation.

The design of TorchDCM follows three principles.
First, the package uses discrete-choice-native abstractions, including choice datasets, utility specifications, beta parameters, nests, random coefficients, thresholds, latent classes, and post-estimation results.
Second, the likelihood kernels are written as vectorized PyTorch computations, allowing a common implementation strategy for ragged choice sets, availability masks, panel likelihoods, simulation draws, and Hessian-based inference.
Third, every implemented estimator is developed together with benchmark cases that compare log likelihoods, parameters, probabilities, covariance matrices, standard errors, willingness-to-pay measures, elasticities, and runtime splits against external reference software.

This paper makes three contributions.
First, we introduce TorchDCM as an open-source PyTorch-first package for discrete choice estimation and inference.
Second, we document the package architecture and implementation choices needed to support both standard and advanced choice models, including multinomial, nested, cross-nested, mixed, WTP-space, ordered, latent class, and hybrid-choice components.
Third, we build a reproducible benchmark suite on public datasets aligned with Biogeme, Apollo, and R package examples, and we report estimator parity and runtime results across model families.
The current benchmark suite establishes full-estimation parity for MNL, nested logit, cross-nested logit, ordered logit, and ordered probit models, and it establishes shared-draw likelihood parity for mixed-logit models.

The rest of the paper is organized as follows.
Section~\ref{sec:related} positions TorchDCM relative to existing choice-modeling software.
Section~\ref{sec:design} describes the package design.
Section~\ref{sec:implementation} presents implementation details for data handling, likelihood computation, inference, and simulation.
Section~\ref{sec:benchmark_protocol} describes the benchmark data and alignment protocol.
Section~\ref{sec:results} reports computational results.
Sections~\ref{sec:discussion} and~\ref{sec:conclusion} discuss current limitations and future work.
